\documentclass[titlepage,a4paper]{jsarticle}
\usepackage{../../../sty/import}% 各種パッケージインポート
\usepackage{../../../sty/title}% タイトルページの変更
\usepackage{../../../sty/answergrid}
%% タイトルページの変数
% レポートタイトル
\title{心理学特論第11回目出席票}
% 提出日
\expdate{\today}
% 科目名
\subject{心理学特論}
% 分野
\class{情報経営システム工学分野}
% 学年
\grade{M1}
% 学籍番号
\mynumber{24336488}
% 記述者
\author{本間三暉}

\begin{document}
% titleページ作成
\maketitle
\section*{keyword\#1}
ヒヤリハット
\section*{keyword\#2}
ハインリッヒの法則
\section*{心理学特論第11回テーマ「失敗の心理学」前編，可能な範囲で簡潔にまとめましょう！}
\section{「失敗から体験的知識を得ることにより成長の糧にすることができるのが良い失敗」であると，授業では定義しました。
自分の体験や他人の経験で「良い失敗」だったと思えることがあれば，それをなるべくわかりやすく説明して下さい。}
私にとって良い失敗だったと思える経験は，就職活動における面接での失敗である．

就職活動を始めた頃の私は，面接で自分の経験をうまく説明することができなかった．
開発アルバイト，サークル活動，個人開発など，話せる経験自体はいくつかあったが，それぞれをどのように企業に伝えればよいのか整理できていなかった．
そのため，質問に対してしどろもどろになったり，話が長くなりすぎたり，自分が本当に伝えたい強みが相手に伝わらなかったりした．
当時は，面接でうまく話せなかったこと自体が単なる失敗のように感じられ，かなり落ち込んだ．

しかし，その失敗を振り返る中で，面接は単に自分の経験をそのまま話す場ではなく，企業が知りたいことに合わせて自分の経験を整理して伝える場であると気づいた．
そこで，面接の内容を振り返り，どの質問にうまく答えられなかったのか，どの経験を話すと自分の強みが伝わりやすいのかを分析するようにした．
また，自分の就職活動の指針や，企業側から送られてきた改善点も確認し，自分の話し方や伝える内容を少しずつ修正した．

その結果，自分が話すべき内容が明確になり，面接でも以前より落ち着いて答えられるようになった．
例えば，開発経験を話す際にも単に「何を作ったか」だけでなく「どのような課題があり，どのように工夫し，そこから何を学んだのか」まで整理して話せるようになった．
その積み重ねによって，面接を通過できる機会も増えた．

この経験から，失敗はそのまま放置すれば単なる失敗で終わるが，原因を振り返り次の行動に活かすことで，自分を成長させる材料になると学んだ．
面接での失敗は当時つらい経験であったが，そこから体験的知識を得て改善につなげることができたため，私にとっては良い失敗だったと考える．


\section{失敗をごまかそうと隠蔽したり，歪曲したり，単純化することで，重大な失敗に発展してしまうこともあります。
現在の社会ニュースや過去の事例から，対応や判断の誤りが顕著だったものを取り上げて，そのしくじりポイントや概要をまとめて下さい。}

私が対応や判断の誤りが顕著だった事例として取り上げるのは，STAP細胞をめぐる研究不正問題である．

STAP細胞とは，体細胞に弱酸性の刺激などを与えることで，多能性を持つ細胞に変化させられるとされた細胞である．
もし本当に実現できれば，再生医療などに大きな影響を与える可能性があるとして，2014年に大きく報道された．
しかし，論文発表後，画像の不自然さ，データの扱い，再現実験の困難さなど，複数の疑義が指摘された．
その後，理化学研究所の調査により一部で研究不正が認定され，最終的にNatureに掲載されたSTAP細胞に関する論文2本は撤回された．

この事例のしくじりポイントは，単に実験がうまくいかなかったことではなく，疑義が出た後の対応にあると考える．
科学研究では，新しい発見が後から否定されたり，再現できなかったりすること自体はあり得る．
本来であれば，その時点でデータや実験手順を整理し，どこまでが確かで，どこからが不確かなのかを明確にする必要がある．
しかしSTAP細胞の問題では，論文中の画像やデータの誤り，実験の再現性，研究チーム内の責任分担などが十分に整理されないまま，大きな成果として発表されてしまった．

また，疑義が出た後も「STAP現象が本当に存在するか」という問題と「論文のデータや画像の扱いに不正や重大な誤りがあったか」という問題が混ざってしまったことも，問題を大きくした要因だと考える．
現象の有無を検証することと，論文作成上の不備や不正を認めて訂正することは，本来は分けて考えるべきである．
しかし，その切り分けが十分にできなかったため，研究そのものだけでなく，研究機関や科学への信頼も大きく損なわれた．

この事例から学べることは，失敗や誤りが起きたときに，それを小さく見せようとしたり，曖昧な説明で済ませたりすると，かえって重大な失敗に発展するということである．
特に科学研究では，結果が正しいかどうかだけでなく，その結果に至る過程が透明であることが重要である．
間違いがあった場合には，早い段階で誤りを認め，原因を調べ，どこまでが事実でどこからが不確実なのかを説明する必要がある．

STAP細胞の問題は，研究上の失敗に加えて，対応や説明の失敗が重なった事例である．
そのため，単なる実験の失敗ではなく，失敗への向き合い方を誤ることで，社会的な信頼の喪失にまでつながった「悪い失敗」の事例だと考える．

\newpage
\section{心理学特論課題サイコロワーク7回答記入欄}
\AnswerGridFilled{
  {動機づけ}
  {自己実現}
  {X理論}
  {Y理論}
  {衛生要因}
  {モラルハザード}
  {結晶性知能}
  {代表性ヒューリスティック}
  {ディレールメント}
  {シャーデンフロイデ}
  {認知的不協和}
  {編集}
  {ハウスマネー効果}
  {事故奉仕バイアス}
  {トリガー戦略}
  {ジョハリの窓}
  {AIDMA}
  {フレームング効果}
  {PPM}
  {イノベーター理論}
}

\end{document}